\documentclass[11pt]{article}

\usepackage[margin=1in]{geometry}
\usepackage{amsmath,amssymb}

\title{Candidate Compact Source Package for Primitive Dirichlet Pairs}
\author{Working note in \texttt{grh/}}
\date{2026-04-24}

\begin{document}

\maketitle

\begin{abstract}
This note records the current theorem-facing shape of the first non-zeta source
target in the project: a compact-interval paired source package for primitive
Dirichlet characters. The note is intentionally scoped as a blueprint, not as a
closed theorem. It isolates the paired quotient object, the compact-interval
source identity one would like to prove, and the exact burdens that remain open.
\end{abstract}

\section{The paired object}

For a primitive Dirichlet character \(\chi\), define
\[
\Xi_\chi(s):=\Lambda(s,\chi)\Lambda(s,\overline\chi),
\qquad
\Phi_\chi^{\mathrm{pair}}(s):=\frac{\Xi_\chi(2s-1)}{\Xi_\chi(2s)}.
\]

This is the cleanest theorem-facing paired object from current scope alone. It
packages the contragredient pair into one scalar self-dual boundary phase.

\section{Blueprint theorem}

The safe endpoint is not yet a closed theorem but a compact-interval source
package blueprint of the following form.

Fix a compact interval \(I\subset\mathbb R\) on which the paired boundary phase
is regular. Then one seeks a source theorem giving
\[
q_\chi^{\mathrm{pair}}(t)=B_\chi^{\mathrm{pair}}(t)+S_\chi^{\mathrm{pair}}(t),
\qquad t\in I,
\]
where:

\begin{enumerate}
\item \(q_\chi^{\mathrm{pair}}\) is the phase derivative attached to
\(\Phi_\chi^{\mathrm{pair}}\);
\item \(B_\chi^{\mathrm{pair}}\) is one unified paired background term;
\item \(S_\chi^{\mathrm{pair}}\) is the background-subtracted paired strip-edge
source scalar.
\end{enumerate}

The theorem-facing burden inside this blueprint is currently dominated by the
compact-interval summation/regularization step, together with unified paired
background identification and multiplicity bookkeeping.

\section{What is safe to say now}

The current record supports the following safe statements.

\begin{enumerate}
\item The paired object \(\Phi_\chi^{\mathrm{pair}}\) is the cleanest current
non-zeta source candidate from current scope alone.
\item Its boundary phase law is cleaner than the single-channel quotient law.
\item Conditional on an exact paired source theorem, the upstairs one-zero
strip-edge kernel and its positivity are automatic.
\item The safe theorem-facing endpoint is still only the compact source-package
blueprint, not a closed theorem.
\end{enumerate}

\section{What this note does not claim}

This note does not claim:

\begin{enumerate}
\item exact local \(S(m)\)-slot realization;
\item local whitening admissibility;
\item remainder dominance;
\item odd/transverse realization;
\item any ERH/GRH consequence.
\end{enumerate}

Those belong to later theorem steps.

\section{Relationship to the paired slot theorem}

The paired source package is the first step only. The later theorem step would
identify the scalar from this source package with the exact manuscript-style
local value-channel slot.

So the order is:

\begin{enumerate}
\item compact paired source-package blueprint;
\item exact paired slot realization;
\item later remainder / odd-package work.
\end{enumerate}

\section{Provenance}

This note is distilled from:

\begin{enumerate}
\item \texttt{grh/20260423-050630-research-team-grh-dirichlet-tau/notes/paired\_source\_package.md};
\item \texttt{grh/20260423-050630-research-team-grh-dirichlet-tau/notes/dirichlet\_paired\_source.md};
\item \texttt{grh/20260423-050630-research-team-grh-dirichlet-tau/notes/paired\_proof\_plan.md};
\item \texttt{grh/20260423-050630-research-team-grh-dirichlet-tau/agents/20260424-061138-paired-source-package-routeA/report.md};
\item \texttt{grh/20260423-050630-research-team-grh-dirichlet-tau/agents/20260424-061138-paired-source-package-routeB/report.md};
\item \texttt{grh/20260423-050630-research-team-grh-dirichlet-tau/agents/20260424-143617-paired-regularization-plan-routeA/report.md};
\item \texttt{grh/20260423-050630-research-team-grh-dirichlet-tau/agents/20260424-143617-paired-regularization-plan-routeB/report.md};
\item \texttt{grh/20260423-050630-research-team-grh-dirichlet-tau/agents/20260424-143617-paired-regularization-plan-verifier/report.md}.
\end{enumerate}

\end{document}
